\documentclass[letterpaper,12pt]{article}
\usepackage{graphicx}
\usepackage{pdflscape}
\usepackage{amsmath}
\usepackage{amsthm}
\usepackage{lscape}
\usepackage{enumerate}
\usepackage{float}
\usepackage{bm}
\usepackage{grffile}
\usepackage[utf8]{inputenc}
\usepackage{relsize}
\usepackage{tablefootnote}
\usepackage{footnote}
\usepackage{sectsty}
\usepackage{xcolor}
\usepackage{caption}
\usepackage{color}
\usepackage{natbib}
\usepackage{multirow}
\usepackage{bbm}
\usepackage[T1]{fontenc}
\usepackage[multiple]{footmisc}
\usepackage{appendix}
\usepackage{times}
\usepackage{relsize}
\usepackage{caption}
\usepackage{newtxtext,newtxmath}
\usepackage{enumitem}
\usepackage{hyperref}

%%%%%%%%%%%%%%%%%%%%%%%%%%%%%%%%%%%%%%%%%%%%%%%%%%%%%%%%%%%%
%%%   Title Page Information
%%%%%%%%%%%%%%%%%%%%%%%%%%%%%%%%%%%%%%%%%%%%%%%%%%%%%%%%%%%%
\title{\vspace{-.5in} ECON712: Presentation of  replication exercise}
\author{Monday, April 20th @ 9AM}
\date{Room: 10-4560}
%%%%%%%%%%%%%%%%%%%%%%%%%%%%%%%%%%%%%%%%%%%%%%%%%%%%%%%%%%%%
%%%   Double and Single Space Commands
%%%%%%%%%%%%%%%%%%%%%%%%%%%%%%%%%%%%%%%%%%%%%%%%%%%%%%%%%%%%
\newlength{\defbaselineskip}
\setlength{\defbaselineskip}{\baselineskip}
\newcommand{\setlinespacing}[1]{\setlength{\baselineskip}{#1 \defbaselineskip}}
\newcommand{\doublespacing}{\setlength{\baselineskip}{1.3 \defbaselineskip}}
\newcommand{\singlespacing}{\setlength{\baselineskip}{1\defbaselineskip}}
\renewcommand{\thesubsection}{\arabic{subsection}}

%%%%%%%%%%%%%%%%%%%%%%%%%%%%%%%%%%%%%%%%%%%%%%%%%%%%%%%%%%%%
%%%   Set Margins Here
%%%%%%%%%%%%%%%%%%%%%%%%%%%%%%%%%%%%%%%%%%%%%%%%%%%%%%%%%%%%
\textwidth=6.5in
\textheight=9in
\oddsidemargin=0.10in
\evensidemargin=0.10in
\topmargin=-0.5in

%%%%%%%%%%%%%%%%%%%%%%%%%%%%%%%%%%%%%%%%%%%%%%%%%%%%%%%%%%%%%
%%%   Actual Document Begins Here
%%%%%%%%%%%%%%%%%%%%%%%%%%%%%%%%%%%%%%%%%%%%%%%%%%%%%%%%%%%%%
\begin{document}
\pagenumbering{arabic}

\maketitle
\vspace{-.3in}

\section*{Instructions}

(\textbf{This part is pasted from the Replication Exercise guidelines}) On Monday, April 20th each student will present their replication exercise to the class. Presentations should be around 15 minutes, followed by 5 minutes of questions. Your presentation should follow the same structure as the written report:

\begin{enumerate}
    \item A brief summary of the paper and its contribution (3--4 minutes).
    \item A clear explanation of the identification strategy (3-- minutes).
    \item A discussion of the main replicated results, with at least one table or figure (4--5 minutes).
    \item Your proposed improvements and overall assessment (3--4 minutes).
\end{enumerate}

\noindent You are required to use slides (\textbf{please e-mail me the slides before the presentation -- the morning of the presentation is fine}). Make sure your tables and figures are clearly readable. You will be evaluated on the clarity of your presentation, the depth of your understanding of the paper, and the quality of your critical assessment.

\section*{Grading Rubric}
As follows is the grading rubric for your presentation
\begin{itemize}

\item ($15\%$) -- for the quality, organization, and clarity of your presentation

\item ($20\%$) -- for your summary of the paper and general discussion of the paper's contribution to the literature

\item ($45\%$) -- for your discussion of the data, methods, results, and identification-related issues, with the following breakdown:
\begin{itemize}
\item Data presentation ($10\%$): data sources, properties of the data set, and relevant summary statistics included in the replication exercise
\item Methods and identification-related issues ($25\%$): identification strategy and related assumptions for causal inference
\item Results ($10\%$): clear discussion and interpretation of results
\end{itemize}

\item ($15\%$) -- for your overall assessment of the paper including suggested improvements

\item ($5\%$) -- for your ability to respond to questions and your overall conclusion, including a brief discussion of whether the paper increased your interest on the topic

\end{itemize}

\end{document}
